% explainer-source-hash: 97967a9d45db4c4a54a04f7b9e4d97567a9e23aa92e39f3cd61168e4e70269c7
% explainer-source-commit: 0de7610f533d267cc97187ea2d75d1ee6ac1aab2
% explainer-generated: 2026-07-16
% Regenerate with /explain-all (see WIRING.md). Do not edit this stamp by hand:
% it is how the toolchain knows whether this document still matches the code.
\documentclass[11pt,a4paper]{article}
\input{preamble}

\title{\textbf{Flask Online Shop}\\[0.3em]\large The Brief}
\author{A short, accurate tour of a Flask + Stripe storefront}
\date{}

\begin{document}
\maketitle
\thispagestyle{empty}

\section{What it is}

A small web store that sells planets as a joke: nine product cards, accounts with hashed
passwords, and a \emph{Buy} button that hands the visitor to Stripe's own payment page. It
is about 360 lines of Python in a single file, five HTML templates, one stylesheet, and
five pinned dependencies. Its purpose is to wire up the parts a real store needs without a
framework that hides them.

\begin{keyidea}{The honest summary, up front}
The project gets the two genuinely dangerous things right: passwords are bcrypt-hashed and
never stored in plain text, and card numbers never touch this application because Stripe
hosts the payment form. It gets login state badly wrong: "logged in" is a flag on a
database row rather than a property of a visitor's session, so \textbf{one person logging
in makes every visitor appear logged in}. I verified this against the running app. The
defect is severe, localised to six lines, and the most interesting thing here to be able to
explain.
\end{keyidea}

\section{Architecture at a glance}

\begin{figure}[H]
\centering
\resizebox{\textwidth}{!}{
\begin{tikzpicture}[node distance=1cm]
  \node[box, minimum width=2.4cm] (browser) {Browser};
  \node[box, right=3.0cm of browser, minimum width=3.4cm, minimum height=3.0cm] (flask)
    {\textbf{Flask app}\\{\scriptsize main.py}\\[2pt]{\scriptsize 7 routes, 1 model}\\{\scriptsize catalog in Python}};
  \node[store, right=3.0cm of flask] (db) {SQLite\\{\scriptsize one user table}};
  \node[extbox, above=1.5cm of db, minimum width=2.8cm] (stripe) {Stripe\\{\scriptsize hosted Checkout}};
  \node[extbox, below=1.4cm of flask, minimum width=3.2cm] (tpl) {Jinja templates\\{\scriptsize + Bootstrap, CSS}};

  \draw[flow] ([yshift=5mm]browser.east) -- node[lbl, above]{request} ([yshift=5mm]flask.west);
  \draw[flow] ([yshift=-5mm]flask.west) -- node[lbl, below]{HTML} ([yshift=-5mm]browser.east);
  \draw[flow] (flask.east) -- node[lbl, above, align=center]{ORM} (db.west);
  \draw[flow] (flask.north east) -- node[lbl, above left, align=center]{create session} (stripe.west);
  \draw[dflow] (browser.north) |- node[lbl, above, pos=0.72, align=center]{buyer redirected; card entered on Stripe's domain} (stripe.north west);
  \draw[flow] (flask.south) -- node[lbl, right]{render} (tpl.north);
\end{tikzpicture}}
\caption{Solid arrows are this app's traffic. The dashed arrow is the buyer leaving for
Stripe, and it is why no card number ever enters this codebase.}
\end{figure}

\begin{table}[H]
\centering\small
\begin{tabular}{@{}llp{7.4cm}@{}}
\toprule
\textbf{Piece} & \textbf{Tool} & \textbf{Role} \\
\midrule
Web framework & Flask 3.0.3 & Routing, templates, redirects. Small by design. \\
Persistence & SQLAlchemy 2.0.30 & One \code{User} table in a SQLite file under \file{instance/}. \\
Passwords & bcrypt 4.1.3 & Salted, deliberately slow hashing. \\
Payments & stripe 9.12.0 & One API call: create a Checkout Session, redirect to it. \\
Front end & Jinja + Bootstrap 5 & One \file{base.html} shell; nine cards from one looped block. \\
\bottomrule
\end{tabular}
\end{table}

\section{How a purchase flows}

\begin{figure}[H]
\centering
\resizebox{\textwidth}{!}{
\begin{tikzpicture}[node distance=0.85cm]
  \node[box, minimum width=2.6cm, minimum height=1.2cm] (b) {Browser\\{\scriptsize clicks Buy}};
  \node[box, right=3.2cm of b, minimum width=2.9cm, minimum height=1.2cm] (app) {Flask looks up\\the amount\\{\scriptsize server-side}};
  \node[extbox, right=3.2cm of app, minimum width=2.5cm, minimum height=1.2cm] (api) {Stripe API};
  \node[extbox, below=2.2cm of b, minimum width=3.0cm, minimum height=1.2cm] (pay) {Stripe hosted page\\{\scriptsize card entered here}};
  \node[box, below=2.2cm of app, minimum width=2.9cm, minimum height=1.2cm] (succ) {\code{/success}\\{\scriptsize static, unverified}};

  \draw[flow] ([yshift=3mm]b.east) -- node[lbl, above, align=center]{POST\\slug only} ([yshift=3mm]app.west);
  \draw[flow] ([yshift=-3mm]app.west) -- node[lbl, below]{303} ([yshift=-3mm]b.east);
  \draw[flow] (app.east) -- node[lbl, above, align=center]{line item} (api.west);
  \draw[dflow] (b.south) -- node[lbl, left, align=center]{buyer pays} (pay.north);
  \draw[dflow] (pay.east) -- node[lbl, above, align=center]{redirect back} (succ.west);
\end{tikzpicture}}
\caption{The client sends only a slug such as \code{earth}; the server owns the price, so a
buyer cannot tamper with the amount. But the return path to \code{/success} is a plain
browser redirect, which the buyer fully controls.}
\end{figure}

One detail worth keeping: the displayed price and the price charged both derive from a
single \code{unit\_amount} in cents. The README records that a template once hardcoded
"\$69,420" for Uranus while Stripe was told \$42,069. Deriving one from the other makes
that bug unrepresentable, and I confirmed all nine prices now agree with their cent
amounts.

\section{The defect worth understanding}

\begin{lstlisting}[language=Python, caption={The whole of the login-state logic.}]
def is_logged_in():
    return bool(db.session.query(exists().where(User.logged_in == 1)).scalar())
\end{lstlisting}

That query asks whether \emph{any} row in the whole table has \code{logged\_in == 1}. It
never learns who is asking: no cookie is read, no session consulted, no identity involved.
A context processor feeds its answer to the nav bar on every page.

\begin{honest}{Verified: one login logs everybody in}
I drove the running app with three independent HTTP clients, each with its own cookie jar.
Before anyone logs in, all three see "Login." After Alice logs in, all three see "Logout,"
including a brand-new client holding zero cookies that never authenticated.

The README calls this "no real session isolation between browser tabs or devices," which
understates it. Tabs and devices are the symptom. The disease is that login is a property
of the \emph{application}, not of a visitor. Related, and also confirmed: \code{/logout}
requires no authentication and clears whichever flagged row comes back first, so a stranger
can log Alice out; and because the flag lives in the database it survives restarts.

The fix is small: put the user id in Flask's \code{session} (a signed cookie, per-visitor
by construction) and delete the \code{logged\_in} column. Flask's \code{session} is already
imported and used once, only to clear stale flash messages. Because the surrounding
structure is right, changing those few lines fixes every page without touching a template.
\end{honest}

\section{Security posture, stated plainly}

\begin{table}[H]
\centering\small
\begin{tabular}{@{}p{4.3cm}p{2.1cm}p{8.2cm}@{}}
\toprule
\textbf{Area} & \textbf{Verdict} & \textbf{Detail} \\
\midrule
Password storage & Correct & bcrypt with per-user salt; plain text never stored. \\
Card data & Correct & Never touches this app; Stripe hosts the form. \\
Price integrity & Correct & Client sends a slug; the server owns the amount. \\
Secrets & Mostly correct & Read from the environment, \file{.env} gitignored. But the signing key falls back to a publicly known default instead of failing loudly. \\
Session handling & \textbf{Broken} & Global flag, not a session. Verified above. \\
Payment confirmation & \textbf{Broken} & \code{/success} returns 200 to someone who never paid. No webhook, and no orders table to record a payment in even if there were one. \\
Auth on checkout & \textbf{Absent} & Buying requires no login at all; the account system is decorative with respect to purchasing. \\
Input validation & Weak & Client-side only. A 60-character password bypasses the form's \code{maxlength="10"}; an unknown planet slug or a missing password field each return HTTP 500. \\
Data model & Weak & \code{email} has no uniqueness constraint (duplicates orphan the second account); the password column is declared \code{String} but stores bytes, forcing an \code{isinstance} guard at login. \\
Tests & Absent & No test suite. \\
\bottomrule
\end{tabular}
\end{table}

\begin{honest}{The setup instructions do not work as written}
The README says to \code{cp .env.example .env} and edit it. Nothing in this project ever
reads a \file{.env} file: there is no \code{python-dotenv} dependency and no
\code{load\_dotenv()} call anywhere in the tree. A reader who follows the instructions gets
the insecure fallback signing key and an empty Stripe key, and finds out only when checkout
fails. The variables must be exported into the shell instead. Relatedly,
\code{app.config['SESSION\_TYPE'] = 'filesystem'} is a Flask-Session setting, and
Flask-Session is neither installed nor imported, so that line does nothing at all.
\end{honest}

\section{What was and was not verified}

I installed the pinned dependencies into a clean virtual environment and drove the running
server. \textbf{Confirmed working:} it boots; \code{GET /} returns 200 with exactly nine
cards; all nine prices match their Stripe amounts; register, login, wrong password, unknown
email and logout all behave as documented; and the per-user password check (which the
README says was once hardcoded to compare against user 1) is genuinely fixed, tested with
two accounts.

\textbf{Confirmed broken:} everything in the posture table marked broken, weak or absent.

\textbf{Not verified:} a real Stripe payment end to end, which needs a test-mode key this
environment does not have. Without a key the checkout route reaches Stripe and fails only
on authentication, which proves the wiring up to that boundary and no further.
\textbf{Inferred, not observed:} Pluto's \code{unit\_amount} of \code{999999999999999}
cents (rendered as \$9,999,999,999,999.99) exceeds Stripe's documented USD maximum of
\code{99999999} cents, so its Buy button would be expected to fail while the other eight
succeed. The amount is read from the source; the rejection is deduced from Stripe's
published limit rather than seen.

\section{Verdict}

\begin{keyidea}{Where the value is}
This is a competent beginner web app that is right about the two things that genuinely
hurt when they are wrong (plain-text passwords, handling card numbers) and wrong about
state management in a textbook way. Its worth in an interview is not "I built a shop." It
is being able to point at a global mutable flag standing in for a session, explain exactly
why it collapses the moment a second visitor arrives, demonstrate it with three cookie
jars, and name the six-line fix; and to say that "redirect to Stripe, then trust the
redirect back" is the classic payment mistake, whose real fix is a webhook plus an orders
table. Flaws named first by the author read as engineering judgment. The same flaws found
first by an interviewer do not.
\end{keyidea}

\end{document}
